\section{Sparse row defects and the full singular list}\label{tra:regular-ranks}

The preceding estimates keep an irreducible dimension and an index within that block. We next keep the index in the full regular representation throughout the induction. This requires controlling the exceptional input rows before forming the grid trace.

The conclusion is qualitatively comparable with a
bounded regular moment, but the proof follows the singular list itself.
Its sparse argument measures the weighted neighborhoods of the
starting tuple, rather than deleting first interactions. Its entropy
argument groups parallel edges at a subexponential threshold. Together
these estimates permit approximation by an operator of prescribed
rank at each balanced split.

Let $s_j(T_N)$ be the singular values of one sweep, in decreasing order and counting all regular multiplicities.
\begin{theorem}\label{tra:regular-ranks:main}
There is an absolute $c>0$ such that
\begin{equation}
                             s_j(T_N)\le j^{-c}\qquad (1\le j\le N!)\label{tra:regular-ranks:eq1}
\end{equation}
for every power of two $N$.
\end{theorem}
\subsection{Sparse input rows and their exceptional set}
\begin{lemma}\label{tra:regular-ranks:sparse}
\label{tra:regular-ranks:sparse-statement}
Fix an absolute \(\delta>0\), small (we can take \(10^{-4}\)). On the representation indexed by a first row of length \(N-k\), \(1\le k\le N^{1-\delta}\), the following estimate holds
\begin{equation}
                          \|T_N\|\le \exp(-b k\log N)\label{tra:regular-ranks:eq2}
\end{equation}
for all sufficiently large \(N\), with fixed \(b>0\). For \(k=1\) this follows from uniform one-card marginals, so we consider \(k\ge2\).
\end{lemma}
\begin{proof}\label{tra:regular-ranks:sparse-proof}
\label{tra:regular-ranks:good-bad-kernel-proof}
Use the representation on \(k\) distinguished cards, that is, on tuples of distinct positions. The representation in \eqref{tra:regular-ranks:eq2} occurs here and does not occur for \(k-1\) cards. Thus we need only work on the space orthogonal to functions dropping at least one coordinate. Write \(x,y\) for input and output tuples, \(P(x,y)\) for the sweep kernel on the tuples and \(P_I\) for the marginals on subtuples indexed by \(I\subset[k]\). The kernel, between these orthogonal-complement spaces with counting measure, can be replaced by
\[
 D(x,y)=\sum_{I\subset[k]}(-1)^{k-|I|} P_I(x_I,y_I)N^{-(k-|I|)} .
\]
Paths in the sweep between specified endpoints are unique. Draw the \(k\) paths and put an interaction between two indices if their paths use any common two-way gate, irrespective of feasibility. If there is an isolated index, \(D=0\): in every term that index if included multiplies path probability independently by \(1/N\).

Here are bounds for the absolute kernel thus obtained. Put \(p=k/N\). For distinct input tuple \(x\) write \(r(x_i,x_j)\in[1,d]\) for the largest coordinate in which the two positions differ. Call \(x\) good if, except at fewer than \(k/20\) of the indices,
\begin{equation}
                       \sum_{j\ne i} 8\, 2^{-r(x_i,x_j)}
                                      \le p^{1/8}.\label{tra:regular-ranks:eq3}
\end{equation}
On a good input row $x$, fix $I\subseteq[k]$ and sample the
nonnegative kernel $P_I N^{-(k-|I|)}$ before restricting the outputs
to be distinct. The labels in $I$ use one common switch field; the
other labels use independent fair bits. Let $\Gamma$ be the event
that their complete path interaction graph has no isolated vertex.
We prove
\[
 \Pr(\Gamma\mid x)\le p^{c_1k},\qquad p=k/N,
\]
uniformly in $I$. Restricting the final tuple to be injective can
only decrease this row sum.

First fix an order in which to reveal the full paths. Given the paths
already exposed, the switch coins used by the exposed labels in $I$
are fixed, and all other
switch coins remain independent and fair. This follows directly from
the path specification: it prescribes a value for each visited switch,
and imposes no condition on any other switch. The exposed independent
walkers prescribe only their own private bits.
Explore the next path $i$ until it first reaches a switch visited by
an exposed path. Before that time, its updated bits are fresh fair
bits, whether $i$ uses the common field or private bits.

For an earlier path $j$, the two paths cannot share a switch before
$r(x_i,x_j)$: a higher unupdated bit still separates them. At time
$t\ge r(x_i,x_j)$, reaching that path's switch requires one
specified prefix of $t-1$ updated bits. The event that this is a
first encounter is a subset of that prefix event in the stopped
exploration. Its conditional probability is therefore at most
$2^{-(t-1)}$. We have not conditioned on avoiding the earlier
switches; that avoidance remains part of the event being bounded.
A union bound gives
\[
 \Pr\{i\text{ meets an earlier path}\mid\text{exposed paths}\}
 \le\sum_{j\text{ earlier}}\sum_{t=r(x_i,x_j)}^d2^{-(t-1)}
 \le\sum_{j\ne i}4\,2^{-r(x_i,x_j)}.
\]
For every index satisfying \eqref{tra:regular-ranks:eq3}, this
is at most $p^{1/8}$. Thus these encounter indicators have a
uniform conditional upper bound in every fixed revelation order.

Now choose the revelation order uniformly, independently of the
sampled paths. In any fixed graph without isolated vertices, each
vertex has an earlier neighbor with probability at least $1/2$.
The number $A$ of vertices with an earlier neighbor therefore has
mean at least $k/2$ and is at most $k$. It follows that
$\Pr\{A\ge k/4\}\ge1/3$. On a good row fewer than $k/20$
indices violate \eqref{tra:regular-ranks:eq3}; hence, with this
probability over the order, at least $k/6$ indices satisfying that
inequality meet an earlier path.

For a fixed order and any prescribed set of $h$ such indices,
successive conditioning in revelation order bounds the probability
that all their encounters occur by $p^{h/8}$. Take
$h=\lceil k/6\rceil$ and sum over the at most $2^k$ possible
sets. Averaging the order and using the preceding $1/3$ bound gives
\[
 \Pr(\Gamma\mid x)
 \le3\cdot2^k p^{k/48}\le p^{k/96}
\]
for sufficiently large $N$, since $p\le N^{-\delta}$.
This proves the required good-row estimate. The backward occupation
estimate below will supply the corresponding column bound on the
exceptional input rows.

For the bad rows we give the corresponding small bound on column sum, even without requiring any interactions. Running any of the \(I\)-processes backwards from the column tuple \(y\), the probability that any specified set of \(h\) distinct input sites are all occupied among the \(k\) paths is at most \(p^h\). Apply Lemma~\ref{grid:occupation} to singleton sets, using
$|I|$ jointly switched cards and $k-|I|$ independent walkers in the
reverse coordinate order. The bound counts all endpoints before
restriction to distinct input tuples, so it remains an upper bound
for the restricted kernel.

We include details that, on distinct input sites, having too many violations of \eqref{tra:regular-ranks:eq3} has probability at most \(p^{c_2 k}\) under this bound. Dyadic boxes of size \(2^r\) consist of vectors agreeing in coordinates above \(r\). Every violating site is in some such box with occupancy \(h\ge 2\) and with
\[
                       h\ \ge\ \frac{p^{1/8}\,2^r}{8d}.
\]
Take maximal such boxes. If a bad row is produced, they are disjoint and together witness total occupancy \(q\ge k/20\). For \(l\le q/2\) boxes with sizes \(a_i=2^{r_i}\), occupancies \(h_i\), sum \(q\), union bound over occupied sites costs at most the sum, with box lists counted unordered, of
\[
                p^q \prod_{i=1}^l (N/a_i)(e a_i/h_i)^{h_i}.
\]
Thus summing ordered specifications of sizes and occupancies we divide by \(l!\). Using \(a_i/h_i\le 8d p^{-1/8}\), the site-cost sum at fixed \(q,l\) is at most
\[
                p^q\, 2^{O(q)} (C d^2 p^{-1/8})^q\, N^l/l!
                 \ \le\ (C' d^2)^q\, p^{7q/8-l}.
\]
In the last bound we used \(N=k/p,\ k^l/l!\le \exp(O(q))\) in this range. This proves the stated bad-row probability with room to spare.

The terms of the absolute bound on \(D\) have row and column sums trivially at most 1 (per term); on good rows with required interactions, or on columns when keeping just the bad rows, we have the small bounds on one of these two sums. Therefore the matrix norm is at most \(2^{k+1}p^{c_3 k}\), with fixed \(c_3>0\), by the row/column test. This proves \eqref{tra:regular-ranks:eq2}.
\end{proof}
\subsection{Conditional coloring entropy}
\begin{proposition}\label{tra:regular-ranks:crossing}
\label{tra:regular-ranks:crossing-statement}
Suppose sites, \(N=mn\), are arranged in \(m\) rows and \(n\) columns, with
\(\sqrt{N}/2\le m,n\le2\sqrt N\).
Let \(K,L\) be the row and column permutation groups. Consider orthogonal convolution projections on these groups, and denote also by \(Q_K,Q_L\) the operators given by them in the regular representation on \(G=S_N\). We require that the projections separately on the subgroups be tensor products over the rows and over the columns. Write
\[
                    R=\operatorname{rank}_{K} Q_K,\qquad
                    S=\operatorname{rank}_{L} Q_L
\]
for nonzero ranks in the regular representations on the subgroups. Then
\begin{equation}
 \operatorname{Tr}_{G}(Q_K Q_L Q_K Q_L)
       \ \le\ R S\,\exp\!\left(O\!\left(N^{.995}
                          +\frac{\log(RS)}{\sqrt{\log N}}\right)\right).\label{tra:regular-ranks:eq4}
\end{equation}
\end{proposition}
\begin{proof}\label{tra:regular-ranks:crossing-proof}
\label{tra:regular-ranks:conditional-coloring-proof}
Let $q_K,q_L$ be the coefficient densities of the two projections
on their own groups. Their squared masses are $R,S$.
Lemma~\ref{lem:coefficient-compatibility}, in its projection case,
gives
\begin{equation}
 \operatorname{Tr}_{G}(Q_K Q_L Q_K Q_L)
       \le R S\,\frac{|G|}{|K||L|}\Pr(\text{compatibility}),\label{tra:regular-ranks:eq5}
\end{equation}
where now \(k,l\) are drawn independently with densities \(|q_K|^2/R,\ |q_L|^2/S\). The \(n\) column permutations in \(l\) are independent with individual densities at most their projection ranks \(S_v\), \(\prod S_v=S\).

For clarity, using row moves followed by column moves to describe compatibility, fix \(k\). It gives a bipartite multigraph from old to new columns, with one edge per row at each column: how the card in that row moves to a new column. Degree is \(m\). The permutations in \(l\) assign colors (new rows) to edges at the new-column vertices, giving each color once there. Compatibility is exactly that each color also appears once at each old-column vertex. Indeed this permits performing the two moves in the opposite order, coloring at old columns and then moving within the new rows. Other composition conventions give the same estimate with harmless inversions.

We will need the following entropy bound. Take any distribution \(Y\) on proper such colorings (each-side requirement), writing \(Y_v\) for the incident permutation at each new-column vertex. Write \(D_v=\log(m!)-H(Y_v)\), using nats. Put \(u=\log m,\ L_*=\exp(\sqrt{u})\), calling edge bundles (parallel sets) bad if their size exceeds \(L_*\); let \(b_*\) be the number of edges in bad bundles. For sufficiently large sizes the total correlation satisfies
\begin{equation}
 \sum_v H(Y_v)-H((Y_v)_v)
   \ \ge\ N-b_*-O(Nm^{-.03})
                     - O(1/\sqrt{u})\sum_v D_v .\label{tra:regular-ranks:eq6}
\end{equation}
We give details to get power-saving error along with the small coefficient of the deficits.

Order the edges within a vertex \(v\) with bundles contiguous, and compare revealing its permutation alone versus revealing after previously exposed vertices in a uniform random order of vertices given by iid \([0,1]\) priorities. For an edge \(e=(v,w)\) let \(a\) be the number of remaining colors when it is revealed locally, \(P\) its conditional law given the earlier local colors. Put \(d_e=D(P\|\mathrm{unif}_a)\) (random from these earlier colors), so \(\sum_{e\text{ at }v}\mathbb E d_e=D_v\). The entropy gap in \eqref{tra:regular-ranks:eq6} sums the conditional entropy drops on knowing the earlier vertices, which equal expected conditional relative entropies to these \(P\)'s. The global conditional law has to avoid colors used by earlier vertices at \(w\); hence the drop is at least the expectation of minus log of the \(P\)-mass available after this exclusion. Condition to evaluate this latter expectation on the entire actual coloring and on priority \(z\) of \(v\), and use Jensen on other priorities. Any color at \(w\) whose edge there is from a different vertex is excluded with probability \(z\), giving at least
\[
\begin{gathered}
\int_0^1 -\log(1-z+z t_e)\,dz=1-F(t_e),\\
t_e=P(\text{actual colors on bundle }(v,w)),
\qquad F(t)=\frac{t\log(1/t)}{1-t}.
\end{gathered}
\]
with continuous endpoint conventions. The cost lost here is at most 1. Discard the bad bundles, and also each good bundle meeting the last $\lceil m^{.8}\rceil$ edges at its vertex. The latter removal costs at most $O(m^{.8}+L_*)$ edges per vertex and leaves at least $m^{.8}$ colors at every edge of each retained bundle. We bound the loss for each other edge in a good bundle as follows.

Call colors heavy for \(P\) if \(P(c)>m^{-.3}\). If heavy mass is \(T\), we have
\begin{equation}
                     T\le C d_e/u+C m^{-.2}.\label{tra:regular-ranks:eq7}
\end{equation}
Indeed \(a\ge m^{.8}\), and one may use
\(D(P\|U)=\sum_i[P_i\log(aP_i)-P_i+1/a]\) with nonnegative summands. The contribution of light colors to \(t_e\) is at most \(L_*m^{-.3}\). If \(T\ge L_*^{-5}\), the bound \(t_e\le T+L_*m^{-.3}\), monotonicity of \(F\), and \eqref{tra:regular-ranks:eq7} show
\begin{equation}
                           F(t_e)\le O(d_e/\sqrt u + m^{-.1}).\label{tra:regular-ranks:eq8}
\end{equation}
For a good bundle, write $\mathcal B$ for its edge set and put
$b=|\mathcal B|\le L_*$. For a current edge $e\in\mathcal B$,
let $\mathcal F_e$ be the local history just before revealing it,
and let $a_e$ be the number of remaining colors. Write $P_e$ for
the current conditional law and $T_e$ for its heavy mass. The heavy set
$H_e=\{c:P_e(c)>m^{-.3}\}$ is determined by $\mathcal F_e$ and
has size at most $m^{.3}$. If the heavy mass is less than $L_*^{-5}$,
the loss is $O(m^{-.1})$ unless some still unrevealed edge
$f\in\mathcal B$ has color in $H_e$; on that event the loss is
at most $C L_*^{-4}$.

Let $P_{f|e}$ be the conditional law of the color of $f$ given
$\mathcal F_e$. The same entropy estimate used above gives
\[
 P_{f|e}(H_e)
 \le C\left(m^{-.1}
       +\frac{D(P_{f|e}\|\mathrm{unif}_{a_e})}{u}\right).
\]
The histories are nested. Conditional entropy decreases between
$\mathcal F_e$ and $\mathcal F_f$, while contiguity of the bundle
gives $0\le a_e-a_f\le L_*$. Since the retained edges have
$a_f\ge m^{.8}$,
\[
 \mathbb E D(P_{f|e}\|\mathrm{unif}_{a_e})
 \le \mathbb E d_f+\log(a_e/a_f)
 \le \mathbb E d_f+C L_*/m^{.8}.
\]
A union bound over the remaining edges and then a sum over current
edges therefore charge each $\mathbb E d_f$ at most $L_*$ times.
The small-heavy-mass losses satisfy
\[
 \begin{split}
 \sum_e\mathbb E[\mathrm{loss}_e;\ T_e<L_*^{-5}]
 &\le O(Nm^{-.05})
   +\frac{C L_*^{-4}}{u}
       \sum_{\mathcal B}
       \sum_{e\in\mathcal B}
       \sum_{\substack{f\in\mathcal B\\ f\text{ not before }e}}
                         \mathbb E d_f\\
 &\le O(Nm^{-.05})
   +\frac{C L_*^{-3}}{u}\sum_v D_v\\
 &\le O(Nm^{-.05})
   +\frac{C L_*^{-2}}{u}\sum_v D_v.
 \end{split}
\]
Here $\mathrm{loss}_e$ denotes the bounded loss $F(t_e)$ after the
light-color error has been included in the displayed additive term.
This proves the asserted small-heavy-mass contribution. Together with
\eqref{tra:regular-ranks:eq8} and the discarded edges, it proves
\eqref{tra:regular-ranks:eq6}.

Consequently, for fixed \(k\), under the possibly weighted column laws used in \eqref{tra:regular-ranks:eq5}, we have
\begin{equation}
         \Pr_l(\text{compatibility}\mid k)
           \le\exp\!\left(-N+b_*+O\left(Nm^{-.03}
                                      +(\log S)/\sqrt u\right)\right).\label{tra:regular-ranks:eq9}
\end{equation}
Indeed take the distribution conditioned on compatibility as \(Y\) in \eqref{tra:regular-ranks:eq6}, if it has positive probability. Minus log of that probability is the relative entropy to the product column laws, equaling the total correlation plus sum of local relative entropies to the respective laws. Each \(D_v\) is bounded by the latter local relative entropy plus \(\log S_v\).

It remains to average the bad-edge factor over \(k\). Under uniform row permutations,
\begin{equation}
                         \mathbb E \exp(c_4\sqrt u\, b_*)=O(1)\label{tra:regular-ranks:eq10}
\end{equation}
for a fixed \(c_4>0\). To verify, we can choose bad cells in the table of old against new columns, each with \(h>L_*\) rows moving accordingly. For total \(q\) such prescriptions, assignment probabilities in the independent row permutations are at most \((e/n)^q\), by the falling factorial bound (or are zero). Lists of cell sizes summing to \(q\) can be counted with \(2^q\); choices of cells with factor at most \((n^2)^{q/L_*}\); choices of rows with factor at most \((em/L_*)^q\). Applying this also for the event \(b_*=q\) proves \eqref{tra:regular-ranks:eq10}. Thus under weighted rows (density bounded by \(R\)), averaging \(\exp b_*\) costs at most \(\exp(O(1+\log R/\sqrt u))\), by Hölder. Finally \(|G|/(|K||L|)\le\exp(N+O(\sqrt N\log N))\) by factorial estimates. Equations \eqref{tra:regular-ranks:eq5}, \eqref{tra:regular-ranks:eq9}, \eqref{tra:regular-ranks:eq10} yield \eqref{tra:regular-ranks:eq4}.
\end{proof}
\subsection{Regular singular-number induction}
\begin{proof}[Proof of Theorem~\ref{tra:regular-ranks:main}]
\label{tra:regular-ranks:regular-singular-induction}
Fix thresholds with room between the preceding power exponents; we use
\[
                        X_N=N^{.998}.
\]
We show \eqref{tra:regular-ranks:eq1} by an induction with a small summable exponent loss. More precisely, there will be exponents \(c_N\) bounded away from zero, \(\le a_0=1/32\) (and chosen sufficiently small absolutely), giving \eqref{tra:regular-ranks:eq1} at each size.

Split the bit list into halves as close as possible, forming the \(m\)-by-\(n\) array of (row, column), column bits those updated first. The sweep factors as independent row sweeps then independent column sweeps; their sizes meet the assumption of \eqref{tra:regular-ranks:eq4}. Work in the regular representation of \(G\). By singular/polar decompositions the singular numbers of \(T_N\) are those of \(E D\), where \(D,E\) are positive semidefinite contractions: the output-side singular magnitude for the first factor, respectively the input-side singular magnitude for the second factor (with the unitary pieces on the outside removed). In particular, the operators \(D\) and \(E\) are induced from tensor products of such operators on the row subgroup and column subgroup, respectively.

Suppose we have child bounds \eqref{tra:regular-ranks:eq1} with common \(c'\le a_0\), the smaller of our exponents available for \(m,n\). To keep cutoffs in this proof independent of the possibly tiny size of \(c'\), first bound \(E^r D^r\) with
\[
                                    r=a_0/c'\ge1 .
\]
We will show, for \(x=\log j\ge X_N\),
\begin{equation}
                  s_j(E^r D^r)
                       \le\exp\left(-a_0(1-O(\epsilon_N)) x\right),
                    \qquad \epsilon_N=(\log N)^{-1/4}.\label{tra:regular-ranks:eq11}
\end{equation}

Decompose \(D,E\) into spectral bands, tensor factor by tensor factor. Here are useful conventions to accommodate repetitions and regular representation projections. Separately for each row or column site-group operator use a band per dyadic range of ending indices of distinct nonzero eigenvalues in its regular list, taking index 1 (the constant eigenvalue) by itself. Eigenvalue 1 is simple in such regular lists: attaining norm 1 in a sweep product of switch-averaging projections requires belonging to their common invariant space, and all cube-edge swaps together generate the symmetric group. If rank in one band is \(t\), every eigenvalue in it is at most \((t/2)^{-c'}\), by the ending-index convention. The projections themselves are convolution operators, including for bands constructed this way. Take tensor combinations of bands over rows for \(D\) and over columns for \(E\), obtaining projections \(Q_K,Q_L\) as in \eqref{tra:regular-ranks:eq4}, of ranks \(R,S\) on the subgroups. Write \(y=\log(RS)\) in this paragraph and the next one (distinct from output tuples earlier). There are at most \(M=\exp(O(\sqrt N\log N))\) pairs of tensor combinations, by e.g. counting dyadic bands. Thus \(E^r D^r\) is the sum of these pieces (ignoring zero eigenvalues). Each piece is, up to contractions on the outside, \(Q_L Q_K\) times a scale of at most
\[
                                W=\exp(-a_0 y+C\sqrt N).
\]
The fourth-power singular sum of \(Q_L Q_K\) on \(G\) is the trace \eqref{tra:regular-ranks:eq4}.

To check \eqref{tra:regular-ranks:eq11}, approximate each piece to error
\[
                             \zeta=\exp(-a_0(1-10\epsilon_N)x)/M
\]
in norm. It can be dropped if \(W\le\zeta\); otherwise by \eqref{tra:regular-ranks:eq4} its necessary rank for this approximation is at most
\[
                       \exp(y+O(N^{.995}+y/\sqrt{\log N}))
                                         (W/\zeta)^4 .
\]
In this second case
\[
                       y\le (1-10\epsilon_N)x
                                      +O(\sqrt N\log N/a_0).
\]
Including the sum over at most $M$ pieces, the logarithm of the
necessary rank is at most
\[
 (1-4a_0)y+4a_0(1-10\epsilon_N)x
 +O\!\left(N^{.995}+\frac{y}{\sqrt{\log N}}+\sqrt N\log N\right)
 \le (1-8\epsilon_N)x.
\]
To obtain the last inequality, use $1-4a_0>0$ and substitute the
upper bound on $y$. The errors are $o(\epsilon_N x)$ because
$x\ge N^{.998}$ and $a_0=1/32$ is fixed. This gives a total rank less than \(j\) and proves \eqref{tra:regular-ranks:eq11}. All large-size conditions here are independent of \(c'\).

Lemma~\ref{lem:power-product-comparison} transfers the powered
estimate to the original positive product: for every $j$,
\[
 \prod_{i=1}^j s_i(ED)
 \le\left(\prod_{i=1}^j s_i(E^rD^r)\right)^{1/r}.
\]
Hence \eqref{tra:regular-ranks:eq11}, geometric averaging over first \(j\) (using trivial contraction on earlier indices), gives for \(x=\log j\ge2X_N\)
\begin{equation}
                          s_j(T_N)
                              \le \exp(-c'(1-C_0\epsilon_N)x)\label{tra:regular-ranks:eq12}
\end{equation}
with an absolute \(C_0\). For example, all but \(O(j e^{-\epsilon_N x/4})\) of the indices up to \(j\) have log index at least \((1-\epsilon_N/2)x\) and are covered by \eqref{tra:regular-ranks:eq11}.

The induction now controls all indices whose logarithm is at least
$2X_N$. We still have to control the beginning of the regular list.
Multiplicity forces a large irreducible block to reach the range just
handled. The remaining small-dimensional blocks have a long first row
or column, where the sparse estimate or annihilation applies. Any singular value on an irrep of dimension \(f\) occurs with multiplicity at least \(f\). Thus irreps with \(f\ge\exp(2X_N)\) can always have their singular values tested by \eqref{tra:regular-ranks:eq12} at least at that threshold, even for values contributing to the start of the regular list. Smaller dimensions come only from shapes near row or column, and \eqref{tra:regular-ranks:eq2} (or vanishing) handles these.

Here are details of these statements about smaller dimensions. For sufficiently large \(N\), if dimension \(<\exp(2X_N)\), either first row or first column has length \(N-k\) with \(k\le N^{.999}\). Indeed if longest length out of both is less than \(N/10\), the hook formula already gives dimension exponentially large in \(N\). Otherwise put a length at least \(N/10\) first as a row, transposing if needed. If at least \(\lfloor N^{.999}\rfloor\) boxes are left in the tail, take a subdiagram retaining exactly that many in the tail and the same first row; its dimension is a lower bound. Writing row length \(l\), new tail size \(a\), the hook formula gives at least
\(\binom{l+a}{a}\exp(-O(a/(l-a)))\): tail hook lengths give at least dimension 1 there, and the first-row hook additions are on the first \(a\) (leftmost) boxes with total addition \(a\), giving the correction displayed relative to \(l!\). This suffices. These arguments use equality of dimensions on transposing. If first column has length \(N-k\) in that range, the switches annihilate the representation: there are no invariants on the matching subgroup of one layer, since the partition does not dominate the content \((2,2,\ldots,2)\), by Young's rule (take \(N\) large here). Shapes with first row \(N-k\), \(1\le k\le N^{.999}\), are covered by \eqref{tra:regular-ranks:eq2}. The constant shape contributes just the singular number at 1.

For \(2\le j\) with \(x=\log j<2X_N\), take the exponent
\(c''=c'(1-C_1\epsilon_N)\), with \(C_1\ge C_0\). Irreps of dimension with log at least \(2X_N\) have norm at most \(\exp(-c'' x)\), by \eqref{tra:regular-ranks:eq12} with their multiplicity (ceiling/real thresholds can equivalently use integer indices). It suffices to count singular values exceeding this test among row shapes of smaller dimensions. Those levels would require by \eqref{tra:regular-ranks:eq2}
\(1\le k < c'' x/(b\log N)\).
The total dimension in the regular list of nonconstant shapes through any such level \(k\) is at most \(N^{4k}\), say, just by occurrence in tuple representations or dimension bounds and counting shapes. Taking our exponents restricted once for all to less than \(b/16\), the number exceeding the test including the trivial value is then less than \(j\) (if the range of levels is empty just use \(1<j\)). So \eqref{tra:regular-ranks:eq1} holds at all indices with \(c''\).

This proves the induction step, with conditions that \(N\) be larger than an absolute cutoff and exponent loss by a factor \(1-C_1\epsilon_N\). We emphasize that no cutoff for the array or spectral-band estimates was dependent on the choice of starting exponent in this induction. For finitely many smaller powers of two there is a common positive starting exponent, as small as required, by simplicity of norm 1 noted above. We can take cutoff already large enough that loss factors used are at least \(1/2\). Lemma~\ref{lem:dyadic-exponents}, with $\gamma=1/4$, bounds the product of the loss factors away from zero along every recursive branch. Thus the induction proves \eqref{tra:regular-ranks:eq1} as claimed.
\end{proof}
\begin{corollary}\label{tra:regular-ranks:mixing}
For an absolute integer $t$, $t$ sweeps meet the total-variation threshold $1/4$.
\end{corollary}
\begin{proof}\label{tra:regular-ranks:mixing-proof}
\label{tra:regular-ranks:fourier-completion}
Apply the converse direction of
Proposition~\ref{e:nonnormal-moment-conversion} to
\eqref{tra:regular-ranks:eq1}, choosing a singular moment with
nonconstant remainder at most $1/4$. The forward direction then
gives total variation at most $1/4$ after an absolute number of
sweeps. A sweep costs $d$ physical steps, and
Lemma~\ref{lem:support} gives the matching-order lower bound.

\end{proof}
